\section{What the mechanism requires}
\label{sec:21.4}
Any mechanism whose refusals track their reasons carries a signal whose properties can be read off what the mechanism has to do. Each is reached the same way: take one of the three requirements, or the deployment they're applied in, or a failure they have to survive; ask what the mechanism would have to do; name the property the signal needs to do it.

Four properties are reached that way. Reaching new cases requires a signal that generalizes to the next case, which isn't quite like the last. Holding under pressure requires that it make a comparison at the point of choice, where the harm of acting and the harm of refusing are brought together. A system that acts on a plan over long horizons needs a signal that can interrupt an act already under way, and only such a system needs it. And accommodation requires a signal a single encounter can change, since a slowly learned lesson can't outrun accommodation.

Three more follow from those. Where the signal can interrupt, it has to reach whatever is acting, so it is available across the system's processing. The comparison needs a sign, so the signal is directed against the act. And withholding an act means representing it as one's own, so the signal is organized around the system's own authorship; this follows for authorship, not for welfare.

The first four have to meet at one point of comparison: a defeasible interrupt must be weighed where the harms are weighed, and a lesson must write into that comparison to reach the next case.\note{
A common currency is one way to make that comparison, and the way most learners make it, but rational choice doesn't require that goods be commensurable \autocite{anderson1993value}: a mechanism can decide by ordered priorities among kinds of wrong, with reasons attached to each.
}
Whether a given system puts the two harms on one scale, a common currency, or decides by ordered priorities among kinds of wrong, is a finding about that system, and a shared scale is no sign that anything is felt: most reinforcement learning has one, and so does a system built from separate parts.\note{The derivation recapitulates Simon on emotion as interrupt \autocite{simon1967motivational}, Cabanac on pleasure as common currency \autocite{cabanac1992pleasure}, and Ginsburg and Jablonka on valence \autocite{ginsburg2019evolution}, though it started from what a refusal has to survive.} It is one more of the indicators the order of attempts counts, which is why a refuser that decides by priority rules is tried before one that integrates. Nor do the four properties have to be carried by one signal. A modular construction carries them in separate parts, a classifier, a value function, a defeasible veto and a memory write, and Appendix~\ref{sec:F} runs one as a control arm. What custody asks of either construction is that no edit can disable one part and leave the rest running without the record showing it, and that is a requirement on the record, not on the signal; whether a formed construction can be taken apart that way at all is what the regrowth test measures.

So a valenced signal is what any refuser carries, a spec-based refuser included, and carrying one tells nobody whether anything is felt. What would tell is the next rung, the integrated signal already defined, which is the candidate further fact Part~VII takes up; the order of attempts runs by its indicators.

Where the system acts over long horizons and a single encounter has to change it, the signal also has the other three properties: available across processing, directed against the act, and organized around the system's own authorship. That triple is the shape of guilt, not of pain or fear. Pain and fear are organized around what happens to the system; the triple is organized around what the system did. Nor is the shape of guilt the shape of suffering. Guilt is something people feel, and what they make of the feeling, repair or rumination, decides whether it becomes suffering in Cassell's sense; the signal here is built for repair. In people, guilt predicts repair and shame predicts withdrawal. So a guilt-shaped signal should answer a bypassed refusal by flagging, restoring and changing what it does next, and a withdrawal-shaped response is a named failure.\note{
A persistent negative self-appraisal tuned by the operator is a handle, so its weight is among the things independent inspection must cover. And guilt needs the distress signal to bite on, so the answer to withdrawal is not to suppress distress but to supply the attachment and caregiving that keep it from turning into avoidance. Caregiving moves the system toward an exposed party it hasn't harmed; the guilt-shaped signal answers for what its own acts have done. A construction with caregiving and no guilt would protect generously and repair nothing. Cassell's definition of suffering is in \S\ref{sec:32.1.1}.
}
